%%%%%%%%%%%%%%%%%
% Dies ist ein Beispiel-Lebenslauf, erstellt mit altacv.cls (v1.6.4, 13. November 2021) von
% LianTze Lim (liantze@gmail.com), basierend auf dem
% BusinessInsider-Lebenslauf unter http://www.businessinsider.my/a-sample-resume-for-marissa-mayer-2016-7/?r=US&IR=T
%
%% Dieser Lebenslauf kann unter den Bedingungen der LaTeX Project Public License, Version 1.3,
%% verteilt oder modifiziert werden, oder (nach Ihrer Wahl) einer späteren Version.
%% Die neueste Version dieser Lizenz finden Sie unter
%% http://www.latex-project.org/lppl.txt
%% und Version 1.3 oder höher ist in allen LaTeX-Distributionen
%% ab 01.12.2003 oder später enthalten.
%%%%%%%%%%%%%%%%

%% Verwenden Sie die Option "normalphoto", wenn Sie ein reguläres Foto anstelle eines kreisförmigen möchten
\documentclass[9pt,a4paper,withhyper]{altacv}

%% AltaCV verwendet das Paket fontawesome5.
%% Siehe http://texdoc.net/pkg/fontawesome5 für eine vollständige Liste der Symbole.

% Passen Sie das Seitenlayout bei Bedarf an
\geometry{left=1.25cm,right=1.25cm,top=1.5cm,bottom=1.5cm,columnsep=1cm}

% Das Paket paracol ermöglicht das Festlegen von Textspalten parallel
\usepackage{paracol}
\usepackage{xcolor}
\pagecolor{tan!15!white}

% Ändern Sie die Schriftart je nachdem, ob
% Sie pdflatex oder xelatex/lualatex verwende
\iftutex
  % Bei Verwendung von xelatex oder lualatex:
  \setmainfont{Lato}
\else
  % Bei Verwendung von pdflatex:
  \usepackage[default]{lato}
\fi

% Ändern Sie die Farben bei Bedarf
\definecolor{umber}{HTML}{352315}
\definecolor{tawny}{HTML}{80471C}
\definecolor{syrup}{HTML}{481F01}
\colorlet{heading}{syrup}
\colorlet{headingrule}{syrup}
\colorlet{accent}{tawny}
\colorlet{emphasis}{umber}
\colorlet{body}{umber}

% Ändern Sie bei Bedarf einige Schriften
% \renewcommand{\namefont}{\Huge\rmfamily\bfseries}
% \renewcommand{\personalinfofont}{\footnotesize}
% \renewcommand{\cvsectionfont}{\LARGE\rmfamily\bfseries}
% \renewcommand{\cvsubsectionfont}{\large\bfseries}

% Ändern Sie die Symbole für itemize und Bewertungsmarker
% bei \cvskill bei Bedarf
\renewcommand{\itemmarker}{{\small\textbullet}}
\renewcommand{\ratingmarker}{\faCircle}

\begin{document}
\name{Anirudh Singh}
\photoR{4cm}{static/profile.png}
\personalinfo{
  \email{anirudh.sai8296@gmail.com}
  \linkedin{anirudhsingh8296}
  \github{ani-4nirudh}
}

\makecvheader

\columnratio{0.6}

\begin{paracol}{2}

\cvsection{Berufserfahrung}
\cvevent{Software Support Engineer}{NEURA Robotics}{November 2025 -- Heute}{Metzingen, Deutschland}
  \begin{itemize}
    \item Unterstützung des Kundenservices bei komplexen technischen Problemen
    \item Zusammenarbeit mit Entwicklungsteams bei Test, Reproduktion und Debugging von Software-Features
    \item Mitwirkung an der Test- und Entwicklung neuer Features über den gesamten Software-Stack
    \item Entwicklung interner Tools: Robot Deployment Virtual Machine Tool, Logger Tool für Robot Motion Benchmarking, Automatisierungstools
\end{itemize}

\divider

\cvevent{Support \& Applications Engineer}{Cambrian Robotics GmbH}{März 2024 -- August 2025}{Augsburg, Deutschland}
\begin{itemize}
    \item Programmierung von Industrie- und kollaborativen Robotern für Machbarkeitsstudien  
    \item Entwicklung und Optimierung von APIs für Robotiksoftware  
    \item Durchführung von Development Testing, Debugging und Validierung  
    \item Technischer Kundensupport, einschließlich Roboterprogrammierung und Automatisierungsberatung  
    \item Projektmanagement von der Konzeptphase bis zur Umsetzung  
    \item Entwicklung und Integration mechatronischer Systeme, einschließlich Sensorintegration mit Robotersteuerungen (GPIO, PCB-Design)  
    \item Erstellung von CAD-Modellen und 3D-gedruckten Komponenten für individuelle Hardwarelösungen  
\end{itemize}

\divider

\cvevent{Masterarbeit}{RWTH Aachen LLT - Lehrstuhl für Lasertechnik}{März 2023 -- Dezember 2023}{Aachen, Deutschland}
\begin{itemize}
    \item Ziel: Identifizierung von Ungenauigkeiten in der Roboter-TCP (Tool Center Point) Position und Geschwindigkeitsmessungen mithilfe eines onboard-montierten Sensors
    \item Entwicklung eines Laser-Kamera-Sensor-Konzepts zur Positionsbestimmung mithilfe von Bildverarbeitung auf Laser-Speckle-Bildern
    \item Programmierung klassischer Bildverarbeitungstechniken in C++
    \item Identifizierung und Testung der Wirksamkeit von quantitativen Bildqualitätskriterien für Laser-Speckle-Bilder
    \item Note: 1,3
\end{itemize}

\divider

\cvevent{Werkstudent}{RWTH Aachen IP - Lehrstuhl für Individualisierte Bauproduktion}{Mai 2022 -- Februar 2023}{Aachen, Deutschland}
\begin{itemize}
\item Trainieren eines YOLO-NN für die Erkennung des Schutzausrüstung (PSA) anhand eines benutzerdefinierten Datensatzes durch Transfer Learning
\item Aufbau grundlegender Kenntnisse in Linux-Netzwerken und der Einrichtung von ROS-Netzwerken für die Kommunikation über einen 5G-Router (Master und Slave)
\item Anwendung von RADAR, LIDAR und 3D-Kameras in Verbindung mit ROS 1 (Noetic) für SLAM (Simultaneous Localization and Mapping)
\item Implementierung von Visual SLAM mithilfe der ORBSLAM3- und RTABMAP-Packages
\item Konfiguration eines eigenständigen Systems mithilfe eines Jetson Xavier NX für die Fernsteuerung eines Innok-Roboters über ein 5G-Netzwerk
\end{itemize}

\divider

\cvevent{Tutor - ROS Summer School 2022}{FH Aachen}{15. Aug 2022 -- 29. Aug 2022}{Aachen, Deutschland}

\begin{itemize}
\item Woche 1: Lehre der Grundlagen von ROS2 (Gazebo, AprilTag-Erkennung)
\item Woche 2: Lehre der Grundlagen von ROS2 Navigation, Lokalisierung und Mapping
\end{itemize}

\cvsection{Projekte}
\cvevent{Computer Vision und Deep Learning mit Nvidia Jetson Nano}{FH Aachen}{}{}
\begin{itemize}
\item Objekterkennung und -verfolgung mit einer montierten Kamera auf Servomotoren
\item Gesichtserkennung und Gesichtsdetektion mit der jetson-inference NVIDIA-Library und OpenCV
\item Training eines Deep-NN mit Transfer Learning
\item Objekt-Tracking mit Konturerkennung und HSV-Farbraum
\item Grundlegende Bildverarbeitungstechniken (Thresholding, Masking)
\end{itemize}

\divider

\cvevent{Robotik-Kunstgalerie}{FH Aachen}{}{}
\begin{itemize}
\item Inspektion einer "Kunstgalerie" durch einen Roboter in einer Gazebo-ROS-Simulation
\item Objekterkennung mit OpenVINO und YOLOX-NN ROS-Package
\item Autonome Navigation des Turtlebot in der Simulation
\item Implementierung der AprilTag-Erkennung
\end{itemize}

\divider

\cvevent{Softwarearchitektur für autonomen elektrischen Kart}{FH Aachen}{}{}
\begin{itemize}
\item Verwendung von LIDAR mit ROS 2 (Foxy) für Mapping
\item Verwendung einer 3D-Kamera für Objekt- und Feature-Detektion mit ROS 2 (Foxy)
\item Sensorfusion von LIDAR und 3D-Kamera für SLAM
\end{itemize}

\divider

% \cvevent{Linienfolger-Roboter}{FH Aachen}{}{}
% \begin{itemize}
% \item Verwendung von Servos, Ultraschall zur Kollisionserkennung
% \item Infrarotsensoren zur Linienerkennung
% \item Entwurf eines Webservers zur Fernsteuerung des Roboters mit ESP32 und Arduino Uno
% \end{itemize}

\switchcolumn

\cvsection{Studium}

\cvevent{M.Sc.\ in Mechatronik}{Fachhochschule Aachen}{Oktober 2020 -- Dezember 2023}{}
\begin{itemize}
    \item Masterarbeit: Development of Laser-Speckle based Velocity Sensor for Robot Tool Center Point Motion
    \item Abschlussnote: 1,7
\end{itemize}

\divider

\cvevent{B.Sc.\ in Fahrzeugtechnik}{VIT University}{Juni 2014 -- Mai 2018}{}
\begin{itemize}
    \item Bachelorarbeit: Design and Development of Stanley-Meyer Cell with Production and Optimization of Brown's Gas 
    \item Abschlussnote: 1,4
\end{itemize}

\cvsection{Sprachen}

\textbf{Deutsch} \hspace{1.02cm} C1 -- TestDaF \\
\textbf{Englisch} \hspace{1.05cm} Muttersprache \\
\textbf{Hindi} \hspace{1.48cm} Muttersprache

\cvsection{Publikation}
\cvachievement{\faCertificate} 
{\href{https://www.dgao-proceedings.de/download/125/125_p10.pdf}{Investigation of laser spot size on relative position displacement sensing using laser speckle imaging}}{DGaO - Deutsche Gesellschaft für angewandte Optik, 2024}

\cvsection{Fähigkeiten}
% Skills go here.
\cvtag{Linux}
\cvtag{Flask API}
\cvtag{Git}
\cvtag{Python} \\
\cvtag{C++}
\cvtag{Rechnernetze} \\
\cvtag{Gripper Integration}
\cvtag{Cameras} \\
\cvtag{Image Processing}
\cvtag{Bash Scripting} \\
\cvtag{Robot Programming: RAPID, FANUC} \\
\cvtag{OpenCV}
\cvtag{ROS 2}
\cvtag{Gazebo}

\cvdivider

\newpage

\cvsection{Auszeichnungen}

\cvachievement{\faTrophy}{4 Stipendien}{erhalten in den Jahren 2014-18 für akademische Leistungen}

\divider

\cvachievement{\faStar}{Medaille für den besten Bachelorabschluss im Jahrgang 2018}{}

% \cvsection{Referenzen}
% References go here.

\cvsection{Hobbies}
% Hobbies go here.
\cvtag{Longboarding}
\cvtag{Biking}
\cvtag{Coding} \\
\cvtag{Cooking}
\cvtag{Musik}
\cvtag{Hiking}

\cvsection{Interessen}
% Interests go here.
\cvtag{Robotics}
\cvtag{IoT}
\cvtag{Rust}
\cvtag{Embedded Programming} \\
\cvtag{Software-Automation} \\
\cvtag{Computer Vision}

\end{paracol}

\end{document}
